
\documentclass[final,ignorenonframetext,compress]{beamer}

\mode<beamer>
{
\usetheme{Bonn}
\pgfdeclareimage[height=0.8cm]{university-logo}{images/Logo_UBo_h24_4c-crop}
\setbeamertemplate{footline}[default]
\logo{\pgfuseimage{university-logo}}
\setbeamertemplate{navigation symbols}{}
}
\mode<handout>{
\usetheme{default}
\usecolortheme{dove}
\setbeamercolor{background canvas}{bg=black!5}
\usepackage{pgfpages}
\pgfpagesuselayout{8 on 1}[a4paper,border shrink=5mm]
}


\DeclareMathOperator*{\mymax}{max}
\DeclareMathOperator*{\argmax}{arg\,max}
\DeclareMathOperator*{\argmin}{arg\,min}
\DeclareMathOperator*{\ForAll}{\bigwedge}
\DeclareMathOperator*{\Exists}{\bigvee}
\newcommand{\W}{\mathcal{W}}

\newcommand*\kreis[2]{\unitlength0.7ex\begin{picture}(2.5,2.5)%


	\color{#2}%
	\put(0.75,0.75){\circle*{3}}%
	\color{black}%
	\put(0.75,0.75){\circle{3}}%
	\put(0.7,0.7){\makebox(0,0){#1}}\end{picture}}
	\newcommand\itplus{\kreis{$+$}{green}}
	\newcommand\itminus{\kreis{$-$}{red}}
	\newcommand\prob[2]{\ensuremath{p_{\mathrm{#1}}\left( #2 \right)}}
	\newcommand\bedprob[3]{\ensuremath{p_{\mathrm{#1}}\left( #2 \left|#3\right. \right)}}

    \newcommand\etal{et.\,al.}

	\usepackage[english]{babel}
	% avant, courier, chancery, times, palatino, bookman, 
	% newcent, utopia, charter
	%\usepackage{times}

	\usepackage[latin1]{inputenc}
	\usepackage[T1]{fontenc}
	\usepackage{multicol}
	\usepackage{mdwtab}
	\usepackage{booktabs}
    \usepackage{tabularx}
	\usepackage{txfonts,textcomp}
	\usepackage{amsmath}
	\usepackage{amsfonts}
    \usepackage{natbib}

	\setlength\parindent{0cm}
	\setlength\parskip{2mm}
	\setbeamerfont{gross}{size=\large}
	\setbeamerfont{quetsch}{size=\tiny}
	\setbeamerfont{klein}{size=\footnotesize}
	\newcommand\w[1]{\mathbf{#1}}

	% example list spacing
	\newcommand\ml{\setlength\itemindent{-11mm}}

    %\newcommand{\newblock}{}


	\title{
Learning Depth-Sensitive Conditional Random Fields for Semantic Segmentation of RGB-D Images
    } 
    \author[short]{Andreas M\"uller}%
	\institute[UBO]{%
    Advisor: Sven Behnke\\
    Institute for Autonomous Intelligent Systems\\
	Rheinische Friedrich-Wilhelms-Universit\"at\\[4mm]
	\includegraphics[width=.2\linewidth]{images/Logo_UBo_h24_4c-crop}\\[3mm]
	Bonn, Germany}


	\begin{document}

	\begin{frame}[plain]
		\titlepage
	\end{frame}

    \section*{Introduction}

	\begin{frame}
		\frametitle{Semantic Segmentation of Structure Classes}
        %FIXME
            \begin{figure}
                \begin{onlyenv}<1-2>
                \includegraphics[width=.4 \textwidth]{images/01124_image}
                \end{onlyenv}
                %\begin{onlyenv}<3>
                %\includegraphics[width=.4 \textwidth]{images/00062_pipeline_sp}
                %\end{onlyenv}
                \begin{visibleenv}<2->
                    \includegraphics[width=.4 \textwidth]{images/01124_gt}
                \end{visibleenv}
            \end{figure}

	\end{frame}

    \section{Structured prediction formulation}

    %\begin{frame}
        %\frametitle{Semantic Segmentation as Structured Prediction}
        %FIXME
        %\begin{figure}
            %\includegraphics<1-2>[width=.6\linewidth]{images/scene_sp_gt}
            %\includegraphics<3->[width=.6\linewidth]{images/scene_sp_graph}
        %\end{figure}
    %\end{frame}

    \begin{frame}
        \frametitle{Structured Prediction}
        Learn prediction function of the form
        \[f(x, w) := \argmax_{y \in \mathcal{Y}}  w^T \psi(x, y) \]
        \begin{visibleenv}<2->
            Pairwise model:
            \[w^T \psi(x, y) = \sum_{(i, j) \in E} w_{i,j} \psi_{i,j}(x, y_i, y_j) + \sum_i w_i \psi_i(x, y_i)\]
        \end{visibleenv}
        \begin{visibleenv}<3->
        Objective
        \end{visibleenv}
        \begin{visibleenv}<4->
            \begin{align*}
                \psi_i(x, y_i) &= f(x_i) \otimes 1_{y_i}\\
                \psi_i(x, y_i, y_j) &= f(x_i, x_j) \otimes 1_{y_i} \otimes 1_{y_j}
            \end{align*}
        \end{visibleenv}
        % FIXME

    \end{frame}

    \section{Unary and Pairwise Features}
    \begin{frame}
        \frametitle{Random Forest Unary Potentials}
        \begin{center}
            \includegraphics[width=.6\linewidth]{images/joerg_rf_features}
        \end{center}
        GPU implementation from \citet{stueckler2013}.\\
        Three trees, max-depth 18.
    \end{frame}

    \begin{frame}
        \frametitle{Absolute height potentials}
            \[
                f_{\text{height}}(x_i) = \textbf{n}_{\text{up}}^T \textbf{n}_{x_i}
            \]
        \begin{visibleenv}<2->
        \begin{center}
            \includegraphics[width=.6\linewidth]{images/height_success}\\
            \vspace{3mm}
            \includegraphics[width=.6\linewidth]{images/height_failure}
        \end{center}
        \end{visibleenv}

    \end{frame}


    \begin{frame}
        \frametitle{Color Contrast}
        \[
        f_\text{color}(x_i, x_j) = \exp\left(-\gamma \lVert c_i - c_j \rVert^2\right).
         \]
        \begin{visibleenv}<2->
            \begin{center}
                \includegraphics[width=.4\linewidth]{images/00062_pipline_feature_pairwise_1}
            \end{center}
        \end{visibleenv}

    \end{frame}

    \begin{frame}
        \frametitle{Edge direction}
        \[
            f_\text{direction}(x_i, x_j) = \sphericalangle(\text{pos}_{x_i} - \text{pos}_{x_j}, [0, 1]^T)
         \]
        \begin{visibleenv}<2->
        \begin{center}
            \includegraphics[width=.4\linewidth]{images/00062_pipline_feature_pairwise_4}
        \end{center}
        \end{visibleenv}
    \end{frame}

    \begin{frame}
        \frametitle{Depth Difference}
        \[
            f_\text{depth}(x_i, x_j) = d_i - d_j / Z
         \]
        \begin{visibleenv}<2->
        \begin{center}
            \includegraphics[width=.4\linewidth]{images/00062_pipline_feature_pairwise_2}
        \end{center}
        \end{visibleenv}
    \end{frame}
    

    \begin{frame}
        \frametitle{Pairwise Normal Orientation}
        \[
            f_\text{normals}(x_i, x_j) = \frac{1}{\pi}\sphericalangle(\mathbf{n}_{x_i}, \mathbf{n}_{x_j})
         \]
        \begin{visibleenv}<2->
        \begin{center}
            \includegraphics[width=.6\linewidth]{images/normal_feature}
        \end{center}
        \end{visibleenv}
    \end{frame}


    \begin{frame}
        \frametitle{Summary of Potentials}
        \begin{align*}
            f(x_i) &= f_\text{RandomForest}(x_i) \oplus f_\text{height}(x_i) \oplus 1\\
            f(x_i, x_j) &= f_\text{color}(x_i, x_j) \oplus f_\text{direction}(x_i, x_j) \oplus f_\text{depth}(x_i, x_j) \oplus f_\text{normals}(x_i, x_j) \oplus 1
        \end{align*}
        
    \end{frame}

    \begin{frame}
        \frametitle{How intractable are loopy models?}
        %FIXME
        \begin{itemize}
            \item In general, inference in loopy models is NP-hard.
            \item Learning relies on inference.
            \item Combining several inference algorithms, we can get exact results in practice. %FIXME reference
        \end{itemize}
    \end{frame}

    \begin{frame}
        \frametitle{Summary}
        \begin{center}
            \includegraphics[width=.9\linewidth]{images/teaser}
        \end{center}
    \end{frame}

    \section{Experiments}
    
    \begin{frame}
        \frametitle{Dataset}
        %TODO
    \end{frame}
    

    \begin{frame}
        \frametitle{Experimental Protocol}
        %TODO
    \end{frame}

    \begin{frame}
        \frametitle{Results}
        \begin{table}[t]
        \begin{tabularx}{\linewidth}{@{\extracolsep{\fill}}lcccccc}
        \toprule
                                & \footnotesize ground        & \footnotesize  structure    & \footnotesize furniture     & \footnotesize props         & \footnotesize class avg   & \footnotesize pixel avg\\
        \cmidrule(r){1-7}

        \footnotesize RF                              &         90.8  &   81.6        & 67.9          & 19.9          &  65.0        &  68.3 \\
        \footnotesize RF + SP                         &         92.5  &   83.3        & \textbf{73.8} & 13.9          &  65.7        &  70.1 \\ 
        \footnotesize RF + SP + SVM                   &         94.4  &   79.1        & 64.2          & \textbf{44.0} &  70.4        &  70.3 \\
        \footnotesize RF + SP + CRF                   & \textbf{94.9} &   78.9        &          71.1 & 42.7          &\textbf{71.9} &  \textbf{72.3} \\
        \cmidrule(r){1-7}
        \footnotesize \citet{SilbermanECCV12}         &         68    &   59          & 70           & 42            &  59.6        & 58.6 \\
        \footnotesize \citet{couprie-iclr-13}         &         87.3  & \textbf{86.1} & 45.3         & 35.5          &  63.5        & 64.5 \\
        \cmidrule(r){1-7}
        \footnotesize \citet{stueckler2013}$^\dagger$ & \textbf{95.6} &   83.0        & \textbf{75.1}& 14.2          &  67.0        & 70.9 \\

        \bottomrule
        \end{tabularx}
        \end{table}
    \end{frame}

    \begin{frame}
        \frametitle{Qualitative Results}
    %{\setlength{\tabcolsep}{1mm} 
    %\begin{tabu} to \linewidth{@{}XXXXX@{}}
        \begin{tabularx}{\linewidth}{@{\extracolsep{\fill}}cccccc}
    \footnotesize Input&
    \footnotesize Random Forest&
    \footnotesize SVM on SP&
    \footnotesize CRF&
    \footnotesize Ground Truth\\

    \includegraphics[width=.17\textwidth]{images/00845_image.png}&%
    \includegraphics[width=.17\linewidth]{images/00845_pixel.png}&%
    \includegraphics[width=.17\linewidth]{images/00845_svm.png}&%
    \includegraphics[width=.17\linewidth]{images/00845_ssvm.png}&%
    \includegraphics[width=.17\linewidth]{images/00845_gt.png}\\

    \includegraphics[width=.17\textwidth]{images/00781_image.png}&%
    \includegraphics[width=.17\linewidth]{images/00781_pixel.png}&%
    \includegraphics[width=.17\linewidth]{images/00781_svm.png}&%
    \includegraphics[width=.17\linewidth]{images/00781_ssvm.png}&%
    \includegraphics[width=.17\linewidth]{images/00781_gt.png}\\

    \includegraphics[width=.17\textwidth]{images/01331_image.png}&%
    \includegraphics[width=.17\linewidth]{images/01331_pixel.png}&%
    \includegraphics[width=.17\linewidth]{images/01331_svm.png}&%
    \includegraphics[width=.17\linewidth]{images/01331_ssvm.png}&%
    \includegraphics[width=.17\linewidth]{images/01331_gt.png}\\

    %\includegraphics[width=.15\textwidth]{images/00118_image.png}&%
    %\includegraphics[width=.15\linewidth]{images/00118_pixel.png}&%
    %\includegraphics[width=.15\linewidth]{images/00118_svm.png}&%
    %\includegraphics[width=.15\linewidth]{images/00118_ssvm.png}&%
    %\includegraphics[width=.15\linewidth]{images/00118_gt.png}

    \end{tabularx}
        \includegraphics[width=\linewidth]{legend.pdf}
    \end{frame}


    \begin{frame}
    \begin{center}
    \includegraphics[width=.45\linewidth]{images/vertical_alignment_png}
    \includegraphics[width=.45\linewidth]{images/depth_difference_png}
    \end{center}
    \end{frame}


    \section*{Structured Machine Learning in Python}
    \begin{frame}
        \frametitle{PyStruct - simple structured prediction}
        \begin{center}
            \includegraphics[width=\linewidth]{images/example_gallery}\\
            \url{http://pystruct.github.io}
        \end{center}
    \end{frame}

    \section*{Summary}
    \begin{frame}
        \frametitle{Take Home}
        \begin{itemize}
            \item High level semantic unaries.
            \item Learn everything.
            \item Exact learning of CRF possible.
            \item Code available: \url{http://github.com/amueller/segmentation}
        \end{itemize}
    \end{frame}
    
    \begin{frame}
        \begin{center}
            Thank you for your attention.

            Questions?
        \end{center}
    \end{frame}
    \begin{frame}<beamer:0>
        \bibliographystyle{plainnat}
        \bibliography{paper}
    \end{frame}

	\end{document}
